\chapter{Bak--Tang--Wiesenfeld Sandpile Model}
\label{ch:sandpile}

\section{Project aim}

This project implements and studies the Bak--Tang--Wiesenfeld (BTW)
sandpile model on a finite $3\times3$ square lattice. The main aim is
to create a correct simulation and use it to examine the long-term
behaviour of a small dissipative sandpile.

The project considers the following questions:

\begin{enumerate}
    \item How does the mean height change from an initially empty lattice?
    \item When does the system reach a statistically stationary state?
    \item What are the distributions of avalanche topples, area,
    boundary loss and length?
    \item How strongly are these avalanche measurements related?
    \item How does the open boundary affect corner, edge and centre sites?
\end{enumerate}

The $3\times3$ lattice is small enough for individual avalanches to be
checked by hand. However, it also has strong finite-size effects.
Therefore, the results will not be presented as evidence of a universal
power law.

\section{Model and method}

The model uses a $3\times3$ lattice with open boundaries. Each site has
an integer height $z_{ij}$ representing its number of grains. The
simulation begins with an empty lattice.

At each time step, one grain is added to a uniformly selected random
site. A site becomes unstable when its height reaches four. An unstable
site topples according to

\begin{equation}
z_{ij} \rightarrow z_{ij}-4.
\end{equation}

One grain is transferred to each nearest neighbour. Grains directed
outside the lattice leave the system. Toppling continues until all site
heights are below four.

For each avalanche, the program records:

\begin{itemize}
    \item \textbf{topples}: the total number of toppling events;
    \item \textbf{area}: the number of distinct sites that topple;
    \item \textbf{loss}: the number of grains lost through the boundary;
    \item \textbf{length}: the maximum Euclidean distance from the
    grain-addition site to a site that topples.
\end{itemize}

The main simulation uses 10,000 grain additions. The first 100 steps are
treated as a burn-in period and are excluded from the stationary
statistics. Random seeds 42 and 43 were used to check whether the main
conclusions depend strongly on a particular random sequence.

\section{Implementation and validation}

The original skeleton code has been developed into a working
\texttt{SandPile} class. The program can create the lattice, add grains,
topple unstable sites, calculate boundary loss and process a connected
avalanche until the grid is stable.

The implementation was tested using hand-calculated examples. These
included:

\begin{itemize}
    \item toppling at the centre, an edge and a corner;
    \item a connected avalanche involving several sites;
    \item repeated toppling of the same site;
    \item two valid toppling orders applied to the same configuration;
    \item a check that every final grid is stable; and
    \item a check of mass balance.
\end{itemize}

For example, after 1,000 grain additions, one test simulation had a final
mass of 16 and a total boundary loss of 984. Therefore,

\begin{equation}
M_{\mathrm{input}}
=
M_{\mathrm{final}}+M_{\mathrm{loss}}
=
16+984
=
1000.
\end{equation}

The two different valid toppling orders also produced the same final
stable configuration.

\section{Current results}

After the initial transient, the mean height fluctuated around a stable
level rather than continuing to increase. The stationary mean heights
were 1.8150 for seed 42 and 1.8155 for seed 43. This close agreement
suggests that the result is not strongly dependent on one random
sequence.

\begin{table}[h!]
\centering
\begin{tabular}{lcc}
\hline
Quantity & Seed 42 & Seed 43 \\
\hline
Stationary mean height & 1.8150 & 1.8155 \\
Height standard deviation & 0.2482 & 0.2536 \\
Recorded stationary avalanches & 3214 & 3201 \\
Avalanche probability & 0.3246 & 0.3233 \\
Mean boundary loss per step & 0.9996 & 1.0004 \\
\hline
\end{tabular}
\caption{Comparison of the two 10,000-step simulations.}
\label{tab:seed-comparison}
\end{table}

The mean boundary loss is approximately one grain per stationary step.
Since one grain is also added per step, the average input and output are
approximately balanced. This provides evidence that the system has
reached a statistically stationary state.

For seed 42, the mean avalanche measurements were 2.525 topples, 2.518
distinct toppled sites, 3.079 lost grains and a length of 0.866. The
largest recorded avalanche had 10 topplings, covered all nine sites,
lost 12 grains and had a maximum length of $2\sqrt{2}$.

The avalanche measurements were strongly positively correlated. The
correlation between topples and area was 0.999, while the correlations
between topples and loss and between topples and length were 0.933 and
0.877 respectively. The very strong topples--area relationship is
partly caused by the small lattice, where most sites topple at most once
during an avalanche.

The open boundary also produced a spatial pattern. For seed 42, the
mean heights were 1.748 at corner sites, 1.851 at edge sites and 1.937
at the centre. Corner sites have the lowest mean height because a corner
toppling loses two grains directly. An edge toppling loses one grain,
while a centre toppling has no direct boundary loss.

The program has produced plots of the mean height, avalanche
distributions, correlations and the spatial mean-height pattern. These
figures will be selected and added to the final version of the report.

\section{Reproducibility and version control}

The simulation has been reorganised into a parameterised function so
that the number of steps, burn-in period and random seed can be changed.
Numerical results are saved in NumPy \texttt{.npz} format. A separate
plotting script can load the saved data and reproduce the figures
without rerunning the complete simulation.

The report template was forked from the course GitHub repository. The
report is currently being edited in the forked project. The updated
LaTeX files, figures, bibliography and compiled PDF will be committed
to the fork before a pull request is submitted to the original course
repository.

\section{Remaining work}

The next tasks are:

The immediate next step is to select the most useful figures and write clear captions for them. The numerical results have already been saved in a machine-readable NumPy file, so the reported values can be checked and reproduced.

The report will then be completed by adding a discussion of statistical stationarity, dissipation at open boundaries, relationships among avalanche measurements and the limitations of the small lattice. Exact state-space enumeration or a modified model will only be attempted if the core report is complete.
